
\section{Primary Aim}

To develop and validate a \textbf{Composite Quality Index (CQI)} that mathematically unifies \textbf{Correctness}, \textbf{Maintainability}, \textbf{Security}, and \textbf{Efficiency} into a single governance framework for AI-assisted development.

\section{Research Questions}

Based on the critical gaps identified in current literature, this research seeks to address the following questions:

\begin{itemize}
    \item \textbf{RQ1: The Metric Integration Problem} How can disparate quality dimensions, such as execution efficiency, structural maintainability, and security risk, be normalized and integrated into a single, mathematically sound composite metric? 
    \item \textbf{RQ2: The Reasoning-Security Paradox} To what extent do enhanced reasoning capabilities in LLMs correlate with the production of more secure and maintainable codebases, and how does this "reasoning budget" affect the overall quality index? 
    \item \textbf{RQ3: The Context-Efficiency Gap} How does the reliance on snippet-level generation rather than "whole-repository" reasoning contribute to structural degradation and resource inefficiencies in AI-generated code? 
    \item \textbf{RQ4: Language-Specific Variability} How significantly do the inherent characteristics of programming languages influence the weighting of the final composite score? 
\end{itemize}

\section{Research Objectives}

To answer the research questions and fulfill the primary aim, the following objectives have been established:

\begin{itemize}
    \item \textbf{To quantify structural degradation\footnote{Long-term decline in the quality and organization of a software project’s architecture} patterns} by analyzing the relationship between the increase in code duplication and the simultaneous decline in refactoring activity within AI-assisted repositories.
    \item \textbf{To identify novel technical debt categories} unique to LLM-integrated systems, such as "prompt debt," "hallucinated dependencies," and "architectural blindness," for inclusion in the maintainability assessment.
    \item \textbf{To evaluate the "Efficiency Paradox"} through empirical benchmarking of runtime performance, memory usage, and energy consumption across diverse programming languages using frameworks like EffiBench and Mercury.
    \item \textbf{To establish a "Security-Weighted" functional metric} that integrates traditional $pass@k$ success rates with vulnerability prevalence (VR) and severity scores (CVSS). 
    \item \textbf{To formulate and validate a unified mathematical model (CQI)} that assigns dynamic weights to these four pillars, providing developers with a single, reliable signal for AI-generated codebase health. 
\end{itemize}